\section{Conclusion}

\subsection*{Summary}

The project set out to build a teaching assistant for DAT255 that is grounded in the course
material and that can produce four distinct study formats: explanations, multiple-choice
quizzes, long-answer practice with feedback, and flashcards. The implemented system meets the
mandatory course requirements: a use case is identified, a decoder-only transformer is
constructed and trained from random initialisation, and an objective evaluation is conducted.
Beyond the mandatory items, a retrieval-augmented generation pipeline is built around the
course book, four model configurations are compared on the same task data, and the result is
deployed as an interactive Gradio application.

The architectural starting point was the nanoGPT/GPT-2 design, modified to incorporate
components that have since become standard in open-weight models: rotary positional
embeddings, RMSNorm, SwiGLU, stochastic depth, and a key/value cache. The tokeniser was
extended with four task-specific special tokens so that one model can serve all four study
activities from a single prompt convention. Training data for those four formats was
synthesised locally with Qwen2.5-7B-Instruct from the chunked course book, sidestepping the
cost and rate limits of API-based generation.

\subsection*{Reflection on results}

The most informative comparison turned out not to be the absolute judge score, but the gap
between perplexity and format compliance. Validation perplexity decreased smoothly throughout
training while JSON validity for the quiz task fluctuated by tens of percentage points
between adjacent checkpoints. This is consistent with the observation in
\cite{zheng2023judging} that intrinsic language-model metrics correlate weakly with output
quality on structured generation tasks; in retrospect, format compliance should have been the
primary checkpoint-selection criterion from the start rather than an addition introduced
mid-project.

The custom transformer is not expected to outperform Qwen2.5-3B on free-form correctness, and
this is not a meaningful failure: a model two orders of magnitude smaller, trained on a
single GPU, was never going to beat a 3\,B model pretrained on a web-scale corpus. The
purpose of training a model from scratch was to satisfy the mandatory criterion and to make
the engineering trade-offs concrete. The interesting deployment finding is the opposite of
the conventional one --- on the structured tasks, the small model that has \emph{seen} the
target format thousands of times is more reliable at producing parseable output than a much
larger model relying on in-context examples. For a quiz tab in production, ``always emits
valid JSON'' is more valuable than ``occasionally writes more elegant prose''.

\subsection*{Limitations and further work}

Three limitations are worth flagging. First, the synthetic training data inherits whatever
hallucinations Qwen produces; manual inspection of a random subset caught a handful of
factually incorrect quiz answers, and a more thorough validation pass --- ideally with
multiple judges --- would tighten the upper bound on correctness. Second, the LLM-as-judge
metric is convenient but not free of bias: GPT-4 has been shown to favour longer, more
verbose answers, which advantages the larger models in the comparison. A small human
evaluation on a fixed subset would calibrate the absolute numbers. Third, the retrieval
component is currently single-pass; a re-ranker on top of the HNSW shortlist, or a small
fine-tune of the BGE encoder on the QA pairs, would likely improve NDCG without changing
the generator side at all.

Natural extensions of the project include multi-turn conversation in the \emph{ask} tab, an
automated pipeline that re-indexes when course material changes, and replacing Qwen with a
quantised open-weight alternative (for example LLaMA-3 8B in GGUF form) once compute permits,
keeping the same retrieval and evaluation infrastructure in place.
